\documentclass[11pt,a4paper]{article}

\usepackage[utf8]{inputenc}
\usepackage[T1]{fontenc}
\usepackage[italian]{babel}
\usepackage{amsmath, amssymb}
\usepackage{geometry}
\usepackage{booktabs}
\usepackage{array}
\usepackage{graphicx}
\usepackage{hyperref}
\usepackage{xcolor}
\usepackage{enumitem}
\usepackage{caption}

\geometry{margin=2.5cm}
\hypersetup{colorlinks=true, linkcolor=blue, urlcolor=blue}

\title{Guida alla presentazione: il filo narrativo completo del\\
  progetto di modellazione predittiva dell'IFC}
\author{Progetto di Applied Statistics --- Note di preparazione alla presentazione}
\date{}

\begin{document}
\maketitle

\begin{abstract}
\noindent
Questo documento non è un report tecnico: è un supporto per parlare.
Ripercorre l'intero progetto nell'ordine in cui va raccontato a
voce, spiega ogni numero che compare sul poster, dichiara
esplicitamente cosa è stato volutamente escluso e perché, e si
chiude con una sezione di domande-e-risposte anticipate, in modo che
nessuna domanda ragionevole del pubblico colga il presentatore
impreparato. Leggere le sezioni 1--7 per provare la presentazione;
leggere la sezione 8 la sera prima per provare la difesa.
\end{abstract}

\section{La tesi in una frase}

\textbf{La fragilità comunale italiana è prevedibile e leggibile da
un piccolo gruppo di indicatori socio-economici, ma una componente
geografica residua mostra che il fenomeno non è puramente
individuale --- è anche spaziale.}

Tutto il resto del progetto sostiene questa frase. Se di tutta la
presentazione dovesse sopravvivere una sola frase, dovrebbe essere
questa.

\section{La domanda di ricerca e perché conta}

L'Italian Fragility Composite Index (IFC) è un indice ufficiale a
livello comunale. Il progetto pone due domande:

\begin{enumerate}[leftmargin=2em]
  \item L'IFC può essere spiegato/predetto statisticamente da
  indicatori socio-economici indipendenti (il panel GRINS V3)?
  \item Se sì, cosa guida davvero la fragilità comunale, e la
  spiegazione è puramente socio-economica o la geografia gioca un
  ruolo indipendente?
\end{enumerate}

Questo conta perché se un piccolo insieme di indicatori leggibili
riesce a ricostruire la maggior parte della varianza di un indice di
fragilità ufficiale, quell'indice può essere monitorato a basso
costo tra una pubblicazione ufficiale e l'altra, e i driver
identificati diventano leve candidate per le politiche pubbliche.

\section{I dati, in un paragrafo}

\textbf{Target}: valori IFC grezzi per il 2019 e il 2021, da
\texttt{IFC\_Final\_Analysis\_Sorted.xlsx}; coincide con i decili
MFI ufficiali nel 95.6\% dei casi, confermando che è un proxy
fedele dell'indice realmente usato dai decisori pubblici.
\textbf{Predittori}: il panel comunale GRINS V3, 549 colonne grezze,
trasformate in 326 covariate candidate (i conteggi normalizzati per
la popolazione più una trasformazione log1p, i conteggi negativi
normalizzati solo per la popolazione, le variabili di dimensione
log-trasformate, i tassi lasciati invariati, tutto winsorizzato al
1°/99° percentile per limitare l'influenza dei comuni anomali).
\textbf{Campione}: 15{.}782 osservazioni, 7{.}886 comuni unici, due
anni (2019, 2021). La cross-validation è sempre fatta dividendo
comuni interi nei fold, mai per riga --- dividere per riga farebbe
finire l'osservazione 2019 di un comune nel fold di training mentre
la sua gemella 2021 finisce nel fold di test, facendo trapelare
informazione e gonfiando l'R$^2$ in modo silenzioso.

\section{La pipeline di modellazione, raccontata come sequenza di
domande}

Ogni passo esiste perché risponde a una domanda lasciata aperta dal
passo precedente. Va raccontata in questo ordine.

\subsection*{Passo 1 --- Quanto si può spiegare in assoluto?}
Una regressione LASSO (CV a 10 fold, $\lambda_{\text{1se}}$) è
eseguita su tutte le 326 covariate candidate. Ne seleziona 191.
R$^2$ out-of-sample $= 0.825$ (CV ripetuta 5$\times$5
panel-aware, 25 fold totali). Questo è il \emph{benchmark}: il
tetto di ciò che un modello lineare può spiegare con questo set di
feature. La prima reazione del prof era corretta: 191 covariate non
si possono discutere una per una in nessuna presentazione.

\subsection*{Passo 2 --- Lo stesso risultato si può raccontare con
meno variabili?}
Una pipeline iterativa riduce le 191 covariate a un set parsimonioso
di 25, esplicito e riproducibile:
\begin{enumerate}[leftmargin=2em]
  \item \textbf{Potatura VIF}: si elimina la covariata con il
  Variance Inflation Factor più alto, si ricalcola, si ripete, finché
  tutti i VIF restanti sono sotto 10 (34 covariate eliminate così).
  \item \textbf{Ranking top-N}: i superstiti vengono ordinati per
  effetto standardizzato $|\beta \cdot \mathrm{SD}|$ e si tengono i
  primi 25 --- è ciò che rende le variabili comparabili tra loro
  nonostante vivano su scale naturali diverse.
  \item \textbf{Filtro di significatività}: ogni covariata che
  risulta non significativa una volta che le altre 24 sono nel
  modello viene automaticamente sostituita dal candidato successivo
  nel ranking, ripetendo finché tutte le 25 sono significative.
\end{enumerate}
Una decisione sostanziale esplicita precede tutti e tre i passi: la
superficie comunale grezza viene esclusa a priori, su richiesta
diretta del prof, perché una volta che ogni altra covariata è
espressa come densità o tasso per-capite, un descrittore geografico
statico come la superficie totale non è di per sé un driver di
fragilità --- è una variabile di scala che il resto della pipeline
controlla già implicitamente.

Risultato: 25 covariate, tutte significative a $p < 10^{-6}$, VIF
massimo $= 6.4$. R$^2$ out-of-sample $= 0.784$ --- una perdita
controllata di soli $0.041$ rispetto al benchmark a 191 covariate,
per una riduzione dell'87\% nella dimensione del modello.

\subsection*{Passo 3 --- La geografia aggiunge qualcosa che le 25
covariate non colgono?}
Un modello lineare misto aggiunge una struttura di intercette
casuali sopra gli stessi 25 effetti fissi:
$$
\text{IFC}_i = \beta_0 + \sum_{j=1}^{25}\beta_j X_{ji} +
u_{r(i)} + v_{p(i)\mid r(i)} + \varepsilon_i,
$$
$$
u_r \sim \mathcal{N}(0,\sigma_R^2), \qquad
v_{p\mid r} \sim \mathcal{N}(0,\sigma_P^2),
$$
con $r$ che indicizza 20 regioni e $p$ che indicizza 107 provincie
nidificate nelle regioni. Sono state stimate e confrontate quattro
varianti (solo regione, regione/provincia nidificate, raggruppamento
basato sulla tassonomia GRINS, e la combinazione di geografia e
tassonomia); il modello geografico nidificato è il migliore dei
quattro. R$^2$ out-of-sample $= 0.818$ --- a $0.007$ dal benchmark a
191 covariate, recuperando quasi tutta l'accuratezza persa
diventando parsimoniosi. La correlazione intra-classe attribuibile
alla geografia è circa il 36\%: più di un terzo della varianza
residua dopo i 25 effetti fissi è spiegata da in quale regione e
provincia si trova un comune, non dal valore di nessuna covariata.

\subsection*{Passo 4 --- Quel segnale geografico è già esaurito, o
manca ancora qualcosa?}
Il Moran's $I$ globale sui residui del modello nidificato, usando una
matrice di pesi spaziali kNN-5 simmetrizzata:
$$ I = 0.234, \quad p < 0.001. $$
È positivo e altamente significativo: \emph{dopo} aver controllato
per 25 covariate socio-economiche e intercette casuali
regione/provincia, i comuni vicini condividono ancora fragilità
inspiegata. La geografia non è del tutto assorbita dai confini
amministrativi. La scomposizione locale LISA di un compagno, sugli
stessi residui e stessi pesi, recupera indipendentemente
$I = 0.236$ --- confermando il numero globale --- e lo localizza:
598 comuni formano cluster significativi "High-High" (sotto-predetti,
circondati da altri comuni sotto-predetti), 554 formano cluster
"Low-Low", e 6{.}418 non mostrano un pattern locale significativo.

Qui va inserito un limite esplicito e onesto: i modelli
autoregressivi spaziali espliciti (SAR/CAR) sono stati tentati come
passo successivo naturale, ma non sono convergiti in un tempo di
calcolo praticabile su questa dimensione del dataset, e sono stati
abbandonati piuttosto che riportati con numeri non attendibili.
Questo va dichiarato apertamente, non nascosto.

\subsection*{Passo 5 --- Il modello è statisticamente affidabile, o
solo un R$^2$ fortunato?}
Una batteria di diagnostiche aggiuntive, riassunta brevemente perché
appartiene al report più che al poster: calibrazione per decile
(Pearson/Spearman/CCC di Lin $= 0.885/0.890/0.879$), un intervallo
di predizione empirico al 95\% di ampiezza 7.79 punti IFC, confronto
AIC/BIC che confirma la preferenza per il modello nidificato, e
diagnostica di normalità dei residui (skewness $0.35$, curtosi in
eccesso $2.31$, Shapiro--Wilk $W = 0.982$) che mostra una deviazione
lieve ma reale dalla normalità, documentata per intero con QQ plot
nel report dedicato all'estensione parsimoniosa.

\subsection*{Passo 6 --- Esiste un modo fondamentalmente diverso di
arrivare alla stessa accuratezza?}
In risposta diretta all'osservazione del prof ``PCR, non PCA'':
\begin{itemize}[leftmargin=2em]
  \item \textbf{PCR sulle 326 feature GRINS}: R$^2$ migliore $=
  0.787$ a $K=100$ componenti --- sotto sia il modello parsimonioso
  sia il benchmark. Le componenti principali sono scelte per
  massimizzare la varianza in $X$, non il potere predittivo su $Y$;
  qui le direzioni non coincidono.
  \item \textbf{PCR sui 12 indicatori IFC grezzi} (le stesse
  componenti usate dalla cascata di PCA descrittiva): R$^2 = 0.987$
  a $K=12$, cioè usando \emph{tutta} l'informazione. Il residuo
  dell'1.3\% è informativo, non rumore: rivela che la costruzione
  ufficiale dell'IFC include uno step non lineare (quasi certamente
  una trasformazione a rango percentile) oltre a un'aggregazione
  lineare pesata delle sue 12 componenti.
  \item \textbf{PLS contro PCR, testa a testa}, sulla stessa matrice
  a 326 feature: PLS (che scegli le componenti per covarianza con
  $Y$, non per sola varianza di $X$) domina PCR di circa 5 punti
  percentuali di R$^2$ a ogni numero di componenti. PLS a $K=75$
  raggiunge R$^2=0.825$, pareggiando quasi esattamente il benchmark
  LASSO a 191 covariate ($0.826$).
\end{itemize}
Conclusione di questo passo: la riduzione dimensionale \emph{può}
pareggiare l'accuratezza di LASSO, ma solo sacrificando
l'interpretabilità --- 75 componenti latenti non hanno un significato
sostanziale individuale, mentre le 191 (o le 25 parsimoniose) di
LASSO sono variabili nominate, con segno, leggibili. LASSO resta il
deliverable principale del progetto per motivi di interpretabilità,
non perché sia più accurato.

\section{I cinque driver sostanziali, spiegati}

Effetti standardizzati ($\beta \cdot \mathrm{SD}$): la variazione di
IFC per un aumento di una deviazione standard del predittore, a
parità del resto del modello.

\begin{center}
\begin{tabular}{@{}cllr@{}}
\toprule
Rango & Variabile & Direzione & $\beta \cdot \mathrm{SD}$ \\
\midrule
1 & Massa contribuenti a reddito medio   & meno fragile & $-1.03$ \\
2 & Densità totale dei contribuenti      & meno fragile & $-0.78$ \\
3 & Massa di reddito dichiarato basso    & più fragile  & $+0.57$ \\
4 & Quota di reddito da pensione         & più fragile  & $+0.56$ \\
5 & Densità di piccole imprese di costruzioni & meno fragile & $-0.53$ \\
\bottomrule
\end{tabular}
\end{center}

Lettura: la fragilità è un equilibrio tra \emph{dinamismo
economico} (una classe media ampia e presente; una base ampia di
contribuenti; un tessuto di piccole imprese di costruzioni) e
\emph{disagio economico} (concentrazione di contribuenti a reddito
basso; forte dipendenza da reddito pensionistico piuttosto che da
lavoro). Nessuno dei primi cinque driver è una variabile
demografica o ambientale --- il segnale più forte è inequivocabilmente
fiscale/economico.

\section{Cosa è stato volutamente escluso dal poster, e perché}

\begin{itemize}[leftmargin=2em]
  \item \textbf{PCA descrittiva su GRINS} e \textbf{clustering
  k-means sulle 12 sub-componenti IFC} (contributo di un compagno):
  sono analisi esplorative/di contesto che confermano che la
  fragilità è strutturata geograficamente \emph{prima} che venga
  stimato alcun modello (es. un confronto con la tassonomia di
  urbanizzazione GRINS, indipendente, mostra l'Italia Interna
  concentrata nei decili 9--10, l'Italia Metropolitana nei decili
  1--2). Motivano gli effetti casuali geografici ma non fanno parte
  della pipeline predittiva, quindi vivono nell'introduzione/contesto,
  non nei risultati principali.
  \item \textbf{QQ plot di normalità dei residui e profilo
  Box--Cox}: diagnostiche reali, riportate per intero nel report di
  estensione parsimoniosa, ma troppo tecniche da leggere alla
  distanza di un poster nei pochi secondi che un visitatore dedica a
  ogni pannello.
  \item \textbf{La tabella completa dei coefficienti a 191 covariate
  e il log completo dell'iterazione VIF}: disponibili in appendice al
  report del benchmark; elencarli sul poster non aggiungerebbe
  informazione che un visitatore potrebbe assorbire sul momento.
\end{itemize}

\section{La conclusione da portare a casa, in tre punti}

\begin{enumerate}[leftmargin=2em]
  \item \textbf{GRINS predice bene l'IFC}: R$^2 \approx 0.82$
  out-of-sample con solo 25 covariate interpretabili più una
  gerarchia geografica --- entro un punto percentuale da un
  benchmark a 191 covariate e da un modello PLS a 75 componenti.
  \item \textbf{La narrazione è socio-economica, non demografica o
  ambientale}: massa di classe media, densità di contribuenti e
  tessuto di piccole imprese riducono la fragilità; concentrazione
  di reddito basso e dipendenza da pensione la aumentano.
  \item \textbf{La geografia conta oltre quanto colgono gli
  indicatori}: una correlazione intra-classe del 36\% più un
  Moran's $I = 0.234$ persistente e confermato localmente, dopo aver
  controllato sia per le covariate sia per la geografia
  amministrativa, mostrano che i modelli spaziali espliciti
  (SAR/CAR) sono il passo successivo naturale, attualmente non
  completato.
\end{enumerate}

\section{Domande anticipate e risposte preparate}

\textbf{D: Perché LASSO e non Ridge o una selezione stepwise
semplice?}\\
LASSO esegue selezione di variabili (azzera esattamente alcuni
coefficienti), cosa che Ridge non fa; questo è essenziale perché
l'obiettivo è un modello interpretabile e sparso, non solo a basso
errore. Un controllo con Elastic Net su $\alpha \in \{0.1, 0.25,
0.5, 0.75, 1.0\}$ ha mostrato che l'errore di CV varia meno
dell'1\% su tutto il range, confermando che la scelta di
$\alpha=1$ (LASSO puro) non è arbitraria né fragile.

\textbf{D: Qual è la differenza tra ``feature'' e ``covariata'' nelle
slide?}\\
Una feature è una variabile candidata creata durante l'ingegnerizzazione
(326 in totale), prima di qualsiasi selezione. Una covariata è una
feature che è effettivamente entrata in un modello di regressione
stimato (191 dopo LASSO, 25 dopo il raffinamento parsimonioso). Le
due parole seguono due fasi diverse della stessa pipeline, non due
oggetti diversi.

\textbf{D: Perché la cross-validation è fatta per comune e non per
riga?}\\
Ogni comune appare due volte (2019 e 2021). Dividere per riga
permetterebbe a un anno di un comune di allenare il modello mentre
l'altro anno dello \emph{stesso} comune lo testa, facendo trapelare
informazione e gonfiando l'R$^2$ in modo otticamente favorevole.
Dividere comuni interi nei fold elimina del tutto questa fuga di
informazione.

\textbf{D: Perché escludere la superficie totale?}\\
Su feedback diretto del prof: una volta che ogni altra covariata è
espressa come densità o tasso per-capite, la superficie è un
descrittore di scala statico, non un driver di fragilità in questa
specifica formulazione --- mantenerla avrebbe reintrodotto esattamente
il tipo di covariata mecanica e non sostanziale che il modello
parsimonioso è pensato per evitare.

\textbf{D: Perché la PCR sui 12 indicatori IFC ufficiali raggiunge
solo R$^2=0.987$ e non 1.0, se sono letteralmente le componenti
dell'indice stesso?}\\
Perché la PCR (e la PCA) possono recuperare solo combinazioni
\emph{lineari} delle variabili di input. Il residuo dell'1.3\%
indica che la costruzione ufficiale dell'IFC include uno step non
lineare --- molto plausibilmente una trasformazione a rango
percentile applicata a un certo punto --- che una proiezione lineare
non può riprodurre del tutto. Questo è stato confermato dopo aver
corretto un problema reale di qualità dei dati (righe
comune--anno duplicate nel foglio ufficiale, che gonfiavano
silenziosamente l'errore apparente all'1.8\% prima della
deduplicazione).

\textbf{D: Perché PLS batte PCR così costantemente?}\\
PCR scegli le componenti per massimizzare la varianza dentro $X$ da
sola; PLS scegli le componenti per massimizzare la covarianza tra
$X$ e $Y$. Quando le direzioni di massima varianza e massima
rilevanza predittiva non coincidono --- cosa che il salto osservato a
bassissimi numeri di componenti nella curva di trade-off della PCR
suggerisce essere il caso qui --- PLS ha un vantaggio strutturale che
cresce con il disallineamento dimensionale.

\textbf{D: Se PLS pareggia l'accuratezza di LASSO, perché LASSO
resta il vostro deliverable principale?}\\
L'accuratezza non è l'unico criterio che conta per il progetto. Una
componente PLS è una combinazione lineare di un numero arbitrario di
variabili originali senza un significato individualmente
interpretabile; un coefficiente LASSO è attaccato a una variabile
nominata, con segno, su cui un decisore pubblico può agire. Il
brief del progetto chiede esplicitamente un benchmark
\emph{interpretabile}, quindi LASSO è preferito a parità di
accuratezza.

\textbf{D: Perché il modello spaziale esplicito (SAR/CAR) non è
stato completato?}\\
È stato tentato seriamente; il modello non è convergiuto in un tempo
di calcolo praticabile su un dataset di questa dimensione
(15{.}782 osservazioni panel con una matrice di pesi spaziali a
livello comunale densa). Piuttosto che riportare coefficienti non
attendibili da una stima non convergente, il tentativo è documentato
come un limite onesto e un passo successivo chiaramente
identificato, supportato dall'evidenza del Moran's $I$/LISA che un
tale modello è di fatto giustificato.

\textbf{D: Cosa significa davvero la correlazione intra-classe del
36\%?}\\
Dopo che le 25 covariate a effetto fisso hanno fatto il loro lavoro,
la varianza \emph{residua} dell'IFC si divide circa al 36\% a ``in
quale regione/provincia si trova un comune'' (un'etichetta, non una
covariata misurata) e al 64\% a tutto il resto (variazione comunale
idiosincratica, rumore di misura, e la componente spazialmente
correlata confermata in seguito da Moran's $I$/LISA).

\textbf{D: Quanto dovremmo fidarci dei primi 5 driver --- potrebbero
essere un artefatto della multicollinearità?}\\
La selezione parsimoniosa impone VIF $<10$ per ogni covariata
mantenuta (il VIF massimo osservato è $6.4$), specificamente per
escludere questo dubbio; tutti i cinque driver principali restano
significativi a $p<10^{-6}$ anche dopo che lo step di sostituzione
iterativa ha rimosso i concorrenti collineari.

\end{document}
